\scanpage{18}
la maximalité est une condition loc., et $O,O'$ sont loc. isomorphes. Donc, les classes \deleted{d'équi} à équivalence \unsure{près} de catégories $A$-préadditives ``strictement Deuringiennes'' correspondent aux classes d'isomorphie d'algèbres d'Azumaya \deleted{\unsure{Deuring}} sur $A$ ($A$ anneau de Dedekind) i.e. aux éléments du groupe de Brauer de $K$, $\operatorname{Br}(\unsure{K|A})$.

$C,C'$ deux catégories de Deuring sur $A$
\[
\uline{\operatorname{Equiv}}(C,C')\qquad\text{catégorie des équivalences }A\text{-linéaires}
\]
\inserted{Si elle est $\ne\varnothing$,} c'est une ``catégorie torseur'' sous la catégorie des équivalences $A$-linéaires de $C$ sur elle-même.

Si $C$ a élément $X$, une \deleted{foncteur} \inserted{équivalence} $A$-linéaire \deleted{\unsure{\varphi}} \inserted{$\varphi:C\mathbin{-}C'$} est \unsure{connue} à isom unique près quand on connaît \inserted{$X'=$} $\varphi(X)$, et l'isom
\[
\varphi_X:O=\operatorname{End}(X)\xrightarrow{\sim}O'=\operatorname{End}(X').
\]
Donc la catégorie envisagée, pour $C=C'$, peut s'identifier à la catégorie des couples $(X',u)$ d'un objet de $C'$ et d'un isom \deleted{\unsure{de $O'$ sur $O$}} $O\simeq\operatorname{End}(X')$.
Identifiant $C$ à $\unsure{\mathbb P|P}(O)$, on trouve la \uline{catégorie des} $\uline{O\text{-bimodules}}$ $P$ tels que $P$ soit projectif de rang $1$ pour ses deux structures.

La composition des \deleted{\unsure{isoms}} \inserted{foncteurs} correspond au produit tensoriel des bimodules, les autos du foncteur identique correspondent aux éléments de $O^{\unsure{\natural|\sharp}*}$ (qui, dans le cas Deuring strict, sont les unités de $A$). [Les \uline{endos} du foncteur identique sont évidemment les éléments de $A^{\unsure{\natural|\sharp}}$.]
\mspage{17}
